\documentclass[11pt,a4paper]{article}
\usepackage[utf8]{inputenc}
\usepackage[T1]{fontenc}
\usepackage{lmodern}
\usepackage[french]{babel}
\usepackage{amsmath,amssymb,mathtools}
\usepackage{icomma}
\usepackage{textcomp}
\usepackage{xcolor}
\usepackage[most]{tcolorbox}
\usepackage{enumitem}
\usepackage[a4paper,margin=2.2cm,headheight=26pt,headsep=10pt]{geometry}
\usepackage{fancyhdr}
\usepackage[hidelinks]{hyperref}
\definecolor{primary}{RGB}{222,49,99}
\definecolor{primarydark}{RGB}{150,22,55}
\tcbset{boxbase/.style={enhanced,breakable,boxrule=0.9pt,arc=2.5pt,
  left=3mm,right=3mm,top=2.3mm,bottom=2.3mm,fonttitle=\bfseries,coltitle=white}}
\newtcolorbox[auto counter]{exo}[1][]{boxbase,colback=primary!4,colframe=primary,
  title={\textsc{Exercice}~\thetcbcounter\ifx\relax#1\relax\relax\else\textmd{ --- \itshape #1}\fi}}
\newcommand{\R}{\mathbb{R}}
\newcommand{\euro}{\texteuro}
\newcommand{\ds}{\displaystyle}
\pagestyle{fancy}\fancyhf{}
\fancyhead[L]{\small\textcolor{primarydark}{\textbf{Thème 4} — Systèmes d'équations}}
\fancyhead[R]{\small\textcolor{primarydark}{\textsc{Moyen}}}
\fancyfoot[C]{\small\thepage}
\renewcommand{\headrulewidth}{0.8pt}
\begin{document}

\begin{center}
{\LARGE\bfseries\color{primarydark} Niveau Moyen — Exercices}\\[2pt]
{\color{primary}\rule{0.45\textwidth}{1.2pt}}\\[3pt]
{\small\itshape Combinaison, fractions et pivot de Gauss. Aucune calculatrice.}
\end{center}
\vspace{1mm}

\begin{exo}[systèmes $2\times2$ par combinaison]
Résous :
\[ \text{a) }\begin{cases} 3x+2y=5\\ 2x-3y=12\end{cases}\qquad\qquad
   \text{b) }\begin{cases} 5x-2y=4\\ 3x+y=9\end{cases} \]
\end{exo}

\begin{exo}[chasser les fractions d'abord]
Multiplie chaque équation pour éliminer les dénominateurs, puis résous :
\[ \text{a) }\begin{cases} \ds\frac{x}{2}+\frac{y}{3}=3\\[4pt] x-y=1\end{cases}\qquad\qquad
   \text{b) }\begin{cases} \ds\frac{x+y}{2}=4\\[4pt] \ds\frac{x}{3}-y=-1\end{cases} \]
\end{exo}

\begin{exo}[pivot de Gauss à 3 inconnues]
Résous par la méthode du pivot :
\[ \begin{cases} x+y+z=6\\ 2x+y-z=1\\ x-y+2z=5\end{cases} \]
\end{exo}

\begin{exo}[mettre un problème en système, puis résoudre]
Deux régions $A$ et $B$ ont, ensemble, vacciné $20$ millions de personnes. Si la région $B$ avait
vacciné $4$~millions de personnes de plus, elle aurait vacciné le \emph{triple} de la région $A$.
Nomme les inconnues, écris le système, puis détermine combien de personnes $B$ a vaccinées.
\end{exo}

\begin{exo}[résoudre, puis exploiter le résultat]
On considère $\ds\begin{cases} 3x+2y=5\\ 2x+3y=10\end{cases}$
\begin{enumerate}[label=\alph*),leftmargin=2em,topsep=1pt,itemsep=1pt]
  \item Résous le système.
  \item Déduis-en la valeur de $x^{2}-y^{2}$ (indice : $x^{2}-y^{2}=(x-y)(x+y)$).
\end{enumerate}
\end{exo}

\end{document}
